% =====================================================================
%  preamble.tex —— 导言区（Lead 独占；译者只读，不得修改）
%  梶浦宏成《三角圏の A∞ 増強について》
%  （RIMS 講究録 1918 巻, 2014, pp.42–58）
%  中文译注本 · XeLaTeX 编译
% =====================================================================
\documentclass[11pt,a4paper]{article}

\usepackage[UTF8]{ctex}
\usepackage{amsmath,amssymb,amsthm,mathtools}
\usepackage{mathrsfs}
\usepackage{bm}
\usepackage{graphicx}
\usepackage{tikz-cd}
\usepackage[a4paper,left=2.4cm,right=2.4cm,top=2.6cm,bottom=2.7cm]{geometry}
\usepackage[dvipsnames]{xcolor}
\usepackage{enumitem}
\usepackage{booktabs}
\usepackage{array}
\usepackage{multicol}
\usepackage{fancyhdr}
\usepackage{caption}
\usepackage{titlesec}
\usepackage[hidelinks,unicode]{hyperref}
\hypersetup{pdftitle={三角范畴的 A∞ 增广（中文译注本）},
  pdfauthor={梶浦宏成 著；中文译注：本工程}}

\captionsetup{font=small,labelsep=quad}
\setlength{\emergencystretch}{3em}
\allowdisplaybreaks[1]
\graphicspath{{figures/}}

% ---------- 节标题：原文「N 标题」→「第 N 节 标题」 ----------
\titleformat{\section}{\Large\bfseries}{第 \thesection{} 节}{1em}{}
\titlespacing*{\section}{0pt}{22pt}{12pt}
\titleformat{\subsection}{\large\bfseries}{\thesubsection}{1em}{}
\titlespacing*{\subsection}{0pt}{16pt}{8pt}

% ---------- 页眉页脚 ----------
\pagestyle{fancy}
\fancyhf{}
\fancyhead[L]{\small\itshape On A-infinity enhancements of triangulated categories}
\fancyhead[R]{\small 三角范畴的 A∞ 增广\quad\thepage}
\renewcommand{\headrulewidth}{0.4pt}

% ---------- 定理类环境 ----------
% 原文体例（实测）：定義2.1、定義2.2、定理2.3、定義3.1、補題3.2、問題4.1、
% 定理4.2、定義5.1、例5.2、例5.3、定理5.4、例6.1、命題6.2、命題6.3、
% 注意6.4、定理7.1、系7.2、系7.3 —— 各类共用一个按节编号的计数器。
\theoremstyle{definition}
\newtheorem{definition}{定义}[section]
\newtheorem{theorem}[definition]{定理}
\newtheorem{proposition}[definition]{命题}
\newtheorem{lemma}[definition]{引理}
\newtheorem{corollary}[definition]{推论}
\newtheorem{example}[definition]{例}
\newtheorem{remark}[definition]{注意}
\newtheorem{problem}[definition]{问题}
\renewcommand{\proofname}{证明}

% ---------- 编号公式 ----------
% 原文 (1)–(10) 为全文连续编号，按出现顺序自动编号即可与原文一一对应。
% 带编号的图表（如 p.55 的 図式(10)）同样是 equation 环境里的 tikz-cd。
\renewcommand{\theequation}{\arabic{equation}}

% ---------- 常用宏 ----------
\newcommand{\R}{\mathbb{R}}
\newcommand{\C}{\mathbb{C}}
\newcommand{\Z}{\mathbb{Z}}
\newcommand{\N}{\mathbb{N}}
\newcommand{\Q}{\mathbb{Q}}
\newcommand{\PP}{\mathbb{P}}
\newcommand{\KK}{\mathbb{K}}
\newcommand{\dd}{\mathrm{d}}
\newcommand{\id}{\mathrm{id}}
\newcommand{\rank}{\operatorname{rank}}
\newcommand{\Hom}{\operatorname{Hom}}
\newcommand{\End}{\operatorname{End}}
\newcommand{\Aut}{\operatorname{Aut}}
\newcommand{\Ext}{\operatorname{Ext}}
\newcommand{\Ob}{\operatorname{Ob}}
\newcommand{\Tr}{\operatorname{Tr}}
\newcommand{\Tw}{\operatorname{Tw}}
\newcommand{\Fuk}{\operatorname{Fuk}}
\newcommand{\coh}{\operatorname{coh}}
\newcommand{\Mod}{\operatorname{Mod}}
\newcommand{\cA}{\mathcal{A}}\newcommand{\cB}{\mathcal{B}}\newcommand{\cC}{\mathcal{C}}
\newcommand{\cD}{\mathcal{D}}\newcommand{\cE}{\mathcal{E}}\newcommand{\cF}{\mathcal{F}}
\newcommand{\cG}{\mathcal{G}}\newcommand{\cH}{\mathcal{H}}\newcommand{\cI}{\mathcal{I}}
\newcommand{\cJ}{\mathcal{J}}\newcommand{\cK}{\mathcal{K}}\newcommand{\cL}{\mathcal{L}}
\newcommand{\cM}{\mathcal{M}}\newcommand{\cN}{\mathcal{N}}\newcommand{\cO}{\mathcal{O}}
\newcommand{\cP}{\mathcal{P}}\newcommand{\cQ}{\mathcal{Q}}\newcommand{\cR}{\mathcal{R}}
\newcommand{\cS}{\mathcal{S}}\newcommand{\cT}{\mathcal{T}}\newcommand{\cU}{\mathcal{U}}
\newcommand{\cV}{\mathcal{V}}\newcommand{\cW}{\mathcal{W}}\newcommand{\cX}{\mathcal{X}}
\newcommand{\cY}{\mathcal{Y}}\newcommand{\cZ}{\mathcal{Z}}
\newcommand{\fA}{\mathfrak{A}}\newcommand{\fB}{\mathfrak{B}}\newcommand{\fC}{\mathfrak{C}}
\newcommand{\fD}{\mathfrak{D}}\newcommand{\fM}{\mathfrak{M}}\newcommand{\fm}{\mathfrak{m}}
\newcommand{\ff}{\mathfrak{f}}
\newcommand{\Inf}{\infty}

% A∞ 的两种写法：\ainf 用于正文（A_\infty），\iinf 用于行内上下标
\newcommand{\ainf}{A_{\infty}}
\newcommand{\Ai}{A_\infty}

% 译者注（就地批注；正文不滥用）
% 用法：\transnote{注文}——命令形式，不要写成 \begin{transnote}...\end{transnote}
\newcommand{\transnote}[1]{\par\medskip\noindent{\small\color{Sepia}\textbf{【译者注】}#1\par}\medskip}

\setlength{\parindent}{2em}
\setcounter{tocdepth}{2}
